%! Inverses of Exponential Functions — Unit 5, Lesson 5.2 — Student Packet Cover Sheet
\documentclass[10pt]{article}
\usepackage{precalculus-article}
\usepackage{precalculus-boxes}
\usepackage{ltablex}
\keepXColumns

\begin{document}

% Full-bleed navy banner
\begin{tikzpicture}[remember picture, overlay]
  \fill[navy] (current page.north west)
    rectangle ([yshift=-0.9in]current page.north east);
\end{tikzpicture}
\vspace{-0.8in}

\begin{center}
  {\color{white}\textbf{\LARGE Precalculus}}\\[4pt]
  {\color{skymid}\large\textbf{Unit 5: Logarithmic Functions}}\\[6pt]
  {\color{white}\textbf{\large Lesson 5.2 \quad Inverses of Exponential Functions}}
\end{center}

\vspace{0.22in}
\namedateperiod
\vspace{0.18in}

\begin{learningtargetbox}
\small
\begin{enumerate}[leftmargin=1.5em, itemsep=5pt]
  \item I can construct representations of the inverse of an exponential function with an
        initial value of 1 (i.e., $g(x) = b^x$), and identify the inverse as $f(x) = \log_b x$.
        \hfill\textit{LO 2.10.A}
  \item I can verify that $f(x) = \log_b x$ and $g(x) = b^x$ are inverse functions using
        composition and by swapping ordered pairs.
\end{enumerate}
\end{learningtargetbox}

\vspace{0.14in}

\begin{tocbox}
\small
\renewcommand{\arraystretch}{1.5}
\begin{tabularx}{\linewidth}{@{} c l X r @{}}
\rowcolor{sky}
\textbf{\#} & \textbf{Component} & \textbf{Description} & \textbf{Score} \\
\midrule
1 & Warm-Up        & Spiral review: evaluating logs and exponentials, inverse point & \blank{1.2cm} \\
2 & Guided Notes   & Additive/multiplicative growth; inverse pairs; composition verification & \blank{1.2cm} \\
3 & Group Activity & Constructing and verifying inverses of exponential functions & \blank{1.2cm} \\
4 & Exit Ticket    & Swapped points, composition check, inverse table & \blank{1.2cm} \\
5 & Homework       & Independent practice with inverses of exponential functions & \blank{1.2cm} \\
\midrule
  & \multicolumn{2}{r}{\textbf{Total}} & \blank{1.2cm} \\
\end{tabularx}
\end{tocbox}

\end{document}
